\section{Related Work}

\subsection{Plan Generation for BDI Agents}

Review the traditional use of manually specified plan libraries as the
procedural knowledge used by BDI agents for means--ends reasoning.

Review planning-based approaches that reduce this engineering burden by
constructing or selecting behaviour when an agent encounters a goal that
cannot be handled directly by its existing plans.

Discuss approaches that invoke a classical planner during BDI execution
to obtain a plan for the current state and goal.

Discuss recent work that incorporates generalized planning into BDI
means--ends reasoning to obtain more general behaviour for goal
achievement.

Identify the main distinction in when the generated behaviour is used:
these approaches primarily employ planning as part of reasoning for a
current goal, whereas GP2PL uses planning knowledge to construct
domain-level procedural knowledge before later goal invocations.

Identify the main distinction in the generated representation:
GP2PL produces parameterised and context-sensitive BDI plan templates,
including reusable subsidiary-goal structure, rather than only a plan
for one grounded problem instance.

Position GP2PL as generating a reusable BDI plan library that remains
available across different objects, states, and future goal requests,
while retaining ordinary BDI plan-library execution.


\subsection{Generalized Planning}

Review generalized-planning approaches that generate policies, programs, or reusable subgoal behaviour across families of planning instances.

Position generalized planning as reusable knowledge for BDI library generation.

Distinguish standalone generalized policies from BDI libraries containing context-sensitive plans and callable subsidiary goals.


\subsection{Temporally Extended Goals}

Review work that gives temporally extended goals explicit semantics in agent programming.

Review planning approaches that compile or search over temporal specifications.

Distinguish complete temporal strategy synthesis from coordinating reusable BDI plans to satisfy an accepted temporal goal.

Position temporal monitoring as a mechanism for controlling execution of the generated BDI library rather than replacing the library with a temporal planner.


\subsection{Natural Language for Temporal Goals}

Review approaches that translate natural-language instructions into temporal specifications.

Position natural-language translation as the goal-input interface rather than the central plan-library-generation problem.

Distinguish specification translation from the subsequent generation, composition, and execution of BDI plans.
